\section*{Введение}
\addcontentsline{toc}{section}{Введение}

Финансовые рынки представляют собой информационно насыщенную среду, в которой цены активов изменяются под влиянием не только прошлой ценовой динамики, но и новостей, ожиданий инвесторов, аналитических комментариев и общего информационного фона. С точки зрения классической финансовой теории сложность прогнозирования связана с тем, что цены быстро включают доступную рынку информацию \cite{fama1970efficient}. Поэтому задача прогноза финансового временного ряда не может сводиться только к выбору алгоритма: необходимо понять, какие именно признаки способны нести информацию о будущей динамике.

В классических постановках задачи прогнозирования в качестве признаков часто используются цены открытия и закрытия, максимумы и минимумы за период, объёмы торгов, доходности, лаговые изменения и технические индикаторы. Эти признаки описывают уже наблюдавшееся поведение рынка. С точки зрения машинного обучения качество прогноза существенно зависит от информативности и отбора входных переменных \cite{guyon2003feature}. Поэтому сравнение разных групп признаков является самостоятельной исследовательской задачей, а не только вспомогательным этапом построения модели.

Текстовые источники могут содержать информацию другого типа: оценки событий, реакцию участников рынка, ожидания относительно будущей динамики, упоминания компаний, продуктов, регуляторных решений или макроэкономических факторов. В финансовой литературе показано, что медиа-содержание, тональность новостей и финансовые словари могут быть связаны с поведением рынка \cite{tetlock2007media,loughran2011liability,kearney2014textual,nassirtoussi2014text}. Отдельная линия исследований рассматривает социальные медиа и поисковые данные как источник коллективных ожиданий и настроений \cite{bollen2011twitter,mao2011sentiment}. Поэтому возникает вопрос, можно ли преобразовать финансовые тексты в числовые признаки и использовать их для повышения качества прогноза.

Развитие языковых моделей сделало эту постановку особенно актуальной. Transformer-архитектура \cite{vaswani2017attention} и BERT-подход к построению двунаправленных контекстных представлений текста \cite{devlin2019bert} позволяют получать не только простые оценки тональности, но и embedding-представления. Такие представления могут сохранять больше информации, чем отдельный sentiment-score, поскольку отражают контекст, лексику и смысловые связи внутри сообщения. Для финансовой области важны доменные модели, в частности FinBERT, адаптированные к финансовым текстам \cite{araci2019finbert,yang2020finbert}. В данной работе FinBERT рассматривается не как самостоятельная модель прогнозирования цены, а как инструмент генерации признаков для последующей модели машинного обучения.

Актуальность исследования связана с необходимостью проверить, обладают ли LLM/FinBERT-generated features добавочной предсказательной силой относительно стандартных рыночных признаков. Иными словами, важно не просто построить модель прогноза, а оценить вклад разных групп признаков: что дают только рыночные признаки, есть ли самостоятельный сигнал в текстовых признаках и улучшается ли прогноз при объединении рыночной и текстовой информации. Отдельная задача работы состоит в том, чтобы проверить, насколько этот эффект устойчив при изменении состава рыночного baseline, источника текстов, актива и горизонта прогноза. Такой подход соответствует feature-centric логике: модели используются как инструмент проверки признаков, а не как самоцель.

Объектом исследования являются признаки, используемые в задачах прогнозирования финансовых временных рядов.

Предметом исследования является влияние признаков, извлечённых из финансовых текстов с помощью FinBERT и производных текстовых представлений, на качество прогноза направления движения финансовых активов при разных источниках данных, активах и горизонтах прогноза.

Цель работы --- теоретически обосновать использование LLM/FinBERT-generated features в задаче прогнозирования финансовых временных рядов и эмпирически оценить их влияние на качество прогноза направления цены относительно рыночных признаков.

Для достижения цели решаются следующие задачи:

\begin{enumerate}
    \item Рассмотреть, какие группы признаков используются в задачах прогнозирования финансовых временных рядов.
    \item Обосновать, почему финансовые тексты могут рассматриваться как источник признаков для прогноза.
    \item Описать различие между рыночными признаками, sentiment-признаками и embedding-признаками.
    \item Обосновать использование FinBERT как инструмента извлечения признаков из финансовых текстов.
    \item Сформировать базовые группы признаков: \texttt{prices-only}, \texttt{FinBERT-only} и \texttt{prices + FinBERT}.
    \item Построить экспериментальный протокол, позволяющий сравнить влияние этих групп признаков на прогноз.
    \item Оценить качество моделей на каждой группе признаков с использованием единого хронологического разбиения данных.
    \item Проверить, дают ли FinBERT-признаки самостоятельный сигнал и добавочную предсказательную силу относительно рыночных признаков.
    \item Проверить, сохраняется ли вывод при расширении рыночного baseline за счёт rolling- и volume-based признаков.
    \item Проверить неоднородность эффекта текстовых признаков на нескольких криптоактивах и горизонтах прогноза.
    \item Указать ограничения, связанные с выбором активов, горизонтов, источников текстов и множественным сравнением групп признаков.
    \item Сформулировать направления дальнейшей проверки устойчивости найденных эффектов.
\end{enumerate}

Методологически работа строится как feature-centric эксперимент. Модели используются не как самостоятельный объект исследования, а как инструмент проверки информативности признаков. Поэтому основное внимание уделяется не максимизации accuracy любой ценой, а сопоставлению качества на разных группах признаков при едином протоколе разбиения, нормализации и оценки.

Эмпирическая часть работы состоит из серии последовательных экспериментов. Основной эксперимент выполнен на данных по акции AAPL: целевая переменная задаётся как бинарный индикатор направления изменения цены закрытия на следующий торговый день. Дополнительный эксперимент с расширенным рыночным baseline проверяет, сохраняется ли вклад FinBERT embeddings после добавления rolling- и volume-based признаков и после снижения размерности. Ещё одна дополнительная проверка расширяет постановку на криптоактивы и датасет \texttt{crypto-news}: для BNB, BTC и SOL рассматриваются горизонты 1, 3 и 7 дней.

Для контроля утечки данных во всех экспериментах используется хронологическая логика разбиения: признаки строятся только из информации, доступной к моменту прогноза, а параметры предварительной обработки оцениваются без использования тестовой подвыборки. Такой протокол важен для финансовых временных рядов, поскольку случайное перемешивание наблюдений или обучение преобразований на всей выборке может привести к завышенной оценке качества.

Структура работы соответствует поставленной цели. В первой части рассматриваются подходы к прогнозированию финансовых временных рядов, роль текстовых данных и теоретическое обоснование LLM-generated features. Во второй части описываются данные, группы признаков, методология сравнения и модели как инструмент проверки признаков. Затем последовательно анализируются основной эксперимент, расширение рыночного baseline и мультиактивная проверка на crypto-news. В заключительной части формулируются интерпретация результатов, ограничения исследования и направления дальнейшей проверки.
